\documentclass[12pt]{article}

\usepackage[utf8]{inputenc}
\usepackage[T1]{fontenc}
\usepackage[portuguese]{babel}
\usepackage{geometry}
\geometry{a4paper, margin=2.5cm}

\usepackage{amsmath}
\usepackage{amssymb}
\usepackage{hyperref}
\usepackage{enumitem}
\usepackage{tcolorbox}
\usepackage{listings}
\usepackage{xcolor}

% Configuração de Listings para Python
\lstset{
    language=Python,
    basicstyle=\ttfamily\small,
    keywordstyle=\color{blue},
    stringstyle=\color{red},
    commentstyle=\color{green!60!black},
    breaklines=true,
    showstringspaces=false
}

\hypersetup{
    colorlinks=true,
    linkcolor=blue,
    filecolor=magenta,      
    urlcolor=blue,
}

\begin{document}

\begin{center}
    \vspace{0.5cm}
    \textbf{INSTITUTO FEDERAL DO CEARÁ -- CAMPUS TAUÁ} \\
    \textbf{Tecnologia em Análise e Desenvolvimento de Sistemas} \\
    \textbf{Disciplina: Ciência de Dados / Análise de Dados} \\
    \textbf{Professor: Reginaldo Fernandes} \\
    \vspace{0.5cm}
    \large{\textbf{Exercício Prático de Fixação: Causalidade e Método Científico (N2)}}
    \vspace{0.3cm}
\end{center}

\hrule
\vspace{0.4cm}

\section*{Contexto: Xarope Fitoterápico e a Asma Infantil no Ceará}
Uma equipe de pesquisadores da área de saúde pública conduziu um Experimento Controlado Aleatório (RCT) no Maciço de Baturité, Ceará, para avaliar a eficácia de um xarope fitoterápico feito à base de \textbf{hortelã-chinesa (malvarisco)} no alívio de sintomas de asma infantil.

Foram recrutadas $30$ crianças com diagnóstico de asma persistente leve. As crianças foram distribuídas aleatoriamente por sorteio duplo-cego em dois grupos:
\begin{itemize}
    \item \textbf{Grupo de Tratamento ($N_{\text{tr}} = 15$):} Recebeu o xarope fitoterápico de malvarisco por 4 semanas.
    \item \textbf{Grupo de Controle ($N_{\text{co}} = 15$):} Recebeu um xarope placebo com aroma e aparência idênticos.
\end{itemize}

Ao final das 4 semanas, os médicos avaliaram a melhoria significativa na capacidade respiratória ($1$ para melhora substancial, $0$ para sem melhora):
\begin{itemize}
    \item \textbf{Resultados no Tratamento:} $11$ crianças melhoraram ($73.3\%$) e $4$ não apresentaram melhoras.
    \item \textbf{Resultados no Controle:} $4$ crianças melhoraram ($26.7\%$) e $11$ não apresentaram melhoras.
\end{itemize}

Os vetores de dados originais são:
\begin{itemize}
    \item \lstinline|resultados_tratamento = [1]*11 + [0]*4|
    \item \lstinline|resultados_controle = [1]*4 + [0]*11|
\end{itemize}

\vspace{0.3cm}
\hrule
\vspace{0.5cm}

\section*{Atividades Práticas}

Escreva um script Python modularizado (no arquivo \texttt{pratica\_causalidade.py}) que realize as seguintes etapas e responda às questões no relatório LaTeX correspondente:

\begin{enumerate}[label=\textbf{Tarefa \arabic*}]
    \item \textbf{Estatística Observada:}
    Determine a proporção de melhoria em cada grupo ($p_{\text{tratamento}}$ e $p_{\text{controle}}$) e calcule a estatística de teste observada, definida pela diferença absoluta entre as proporções dos grupos:
    $$\text{Observed\_Diff} = |p_{\text{tratamento}} - p_{\text{controle}}|$$

    \item \textbf{O Modelo de Resultados Potenciais:}
    Utilizando a metáfora do ``bilhete de dois lados'' explicada na Aula 12, descreva o que representa cada lado do bilhete para uma criança que participou do estudo. Explique o conceito de contrafactual.

    \item \textbf{Formulação de Hipóteses:}
    Formule a Hipótese Nula ($H_0$) e a Hipótese Alternativa ($H_1$) para o experimento com o xarope de malvarisco.
    
    \item \textbf{Simulação do Teste de Permutação:}
    Implemente uma simulação com $10.000$ permutações em Python para construir a distribuição nula da estatística de teste (diferença absoluta). Colete as diferenças absolutas geradas por puro acaso.
    
    \item \textbf{Cálculo do $P$-valor e Visualização:}
    \begin{enumerate}[label=\alph*)]
        \item Calcule o $p$-valor empírico bilateral.
        \item Gere um histograma da distribuição nula. Destaque em dourado (\texttt{\#ff9900}) a barra que contém a diferença observada e todas as barras à sua direita. Insira a linha tracejada vermelha no valor observado.
        \item Salve o gráfico como \texttt{histograma\_causalidade.png}.
    \end{enumerate}

    \item \textbf{Inferência de Causalidade:}
    Considerando o nível de significância de $\alpha = 5\%$, qual a decisão estatística? Podemos afirmar que o xarope de malvarisco \textbf{causou} a melhoria na respiração das crianças? Justifique detalhadamente por que a randomização (RCT) nos permite fazer essa afirmação de causalidade, ao contrário de um estudo observacional.
\end{enumerate}

\pagebreak

\section*{Gabarito Comentado e Código de Referência}

\begin{tcolorbox}[colback=yellow!10!white,colframe=yellow!50!black,title=\textbf{Gabarito das Tarefas Teóricas}]
\textbf{Resultados Potenciais (Tarefa 2):}
Cada bilhete de dois lados representa as duas respostas possíveis de uma criança:
\begin{itemize}
    \item \textbf{Lado Esquerdo:} Melhora respiratória caso a criança tome o xarope fitoterápico.
    \item \textbf{Lado Direito:} Melhora respiratória caso a criança tome o placebo.
    \item \textbf{Contrafactual:} É o resultado não observado. Para uma criança sorteada no tratamento, o lado esquerdo é o resultado observado e a resposta sob controle (lado direito) é o contrafactual (desconhecido).
\end{itemize}

\textbf{Decisão e Causalidade (Tarefa 6):}
\begin{itemize}
    \item Diferença Observada: $|0.7333 - 0.2667| = 0.4667$.
    \item P-valor Empírico: $0.0247$ ($2.47\%$).
    \item \textbf{Decisão:} Rejeitamos $H_0$ pois $2.47\% < 5\%$.
    \item \textbf{Causalidade:} Sim, podemos inferir causalidade. Como os pacientes foram alocados aleatoriamente nos dois grupos, a presença de variáveis de confundimento (como severidade inicial da asma, alimentação ou hábitos familiares) foi balanceada igualmente entre os dois grupos. Logo, a única diferença sistemática entre os grupos foi a administração do xarope de malvarisco, provando causa e efeito.
\end{itemize}
\end{tcolorbox}

\subsection*{Script Python de Referência}
\begin{lstlisting}[language=Python]
import numpy as np
import matplotlib.pyplot as plots

# Estilo visual
plots.style.use('seaborn-v0_8-colorblind')
plots.rcParams.update({
    'figure.figsize': (10, 6),
    'axes.labelsize': 12,
    'axes.titlesize': 14,
    'lines.linewidth': 2.5
})

def despine(ax):
    ax.spines['top'].set_visible(False)
    ax.spines['right'].set_visible(False)

# Dados (1 = melhora, 0 = sem melhora)
res_tratamento = np.array([1]*11 + [0]*4)
res_controle = np.array([1]*4 + [0]*11)

# Tarefa 1: Estatistica observada
p_tr = np.mean(res_tratamento)
p_co = np.mean(res_controle)
obs_diff = np.abs(p_tr - p_co)
print(f"Taxa Tratamento: {p_tr:.4f}")
print(f"Taxa Controle: {p_co:.4f}")
print(f"Diferenca Absoluta Observada: {obs_diff:.4f}")

# Tarefa 4: Simulacao de Permutacao
todos_resultados = np.concatenate([res_tratamento, res_controle])
N_tr = len(res_tratamento)
repeticoes = 10000

np.random.seed(42)
distancias_simuladas = np.zeros(repeticoes)

for i in range(repeticoes):
    embaralhado = np.random.permutation(todos_resultados)
    sim_tr = embaralhado[:N_tr]
    sim_co = embaralhado[N_tr:]
    distancias_simuladas[i] = np.abs(np.mean(sim_tr) - np.mean(sim_co))

# Tarefa 5: P-valor e Histograma
p_valor = np.count_nonzero(distancias_simuladas >= obs_diff) / repeticoes
print(f"P-valor empirico: {p_valor:.5f}")

fig, ax = plots.subplots()
bins = np.arange(0.0, 0.7, 0.0667) # bins baseados no tamanho amostral
n, bins_out, patches = ax.hist(distancias_simuladas, bins=bins, edgecolor='white', color='#2c3e50', density=True, label='Sob a Hipotese Nula')

# Destacar a cauda
for patch in patches:
    if (patch.get_x() + patch.get_width()) > obs_diff:
        patch.set_facecolor('#ff9900')

ax.axvline(obs_diff, color='#e74c3c', linestyle='dashed', linewidth=3, label=f'Diferenca Obs. ({obs_diff:.3f})')

from matplotlib.patches import Patch
legenda_patches = [
    Patch(facecolor='#2c3e50', edgecolor='white', label='Valores sob H0'),
    Patch(facecolor='#ff9900', label='Diferenca >= Observada'),
    ax.lines[0]
]
ax.legend(handles=legenda_patches)
ax.set_title('Distribuicao Nula vs. Observado (RCT Hortela-Chinesa)')
ax.set_xlabel('Distancia entre Proporcoes (Valor Absoluto)')
ax.set_ylabel('Densidade')

despine(ax)
plots.tight_layout()
plots.savefig('histograma_causalidade.png', dpi=150)
plots.close()
print("Grafico histograma_causalidade.png salvo!")
\end{lstlisting}

\end{document}
